\section{Preliminaries}\label{sec:prelim}

This section fixes the objects of the paper and the facts about them that the proof uses: walks, hops and walk DAGs
(\cref{sec:prelim-walks}); cuts, labels, $\TCP$ and $\SR$, and the probabilistic method for label cuts
(\cref{sec:prelim-tcp}); owners and prominence (\cref{sec:prelim-prom}); and bands and the random grid
(\cref{sec:prelim-grid}). We follow the eighth paper of the series~\cite{EighthPaper2026} and reuse its formal definitions
and results; its Lean declarations carry the namespace \texttt{BoundedNesting}. All other Lean names are relative to the
namespace \texttt{BoundedStackDepth} of this paper's library.

The proof is an averaging argument. Its random choices are finitely many and equally likely, so every probability is a
number of choices divided by the number of all choices, and every expectation is an average. The Lean statements are the
corresponding inequalities between sums over finite sets; we state them in this form and read them as probabilities.

\subsection{Walks, hops and walk DAGs}\label{sec:prelim-walks}

A \emph{walk} is a map $x\colon\N\to\Z$ with $\abs{x(t+1)-x(t)}\le1$ for every $t$ (\lean{BoundedNesting.IsWalk}); a
step with $x(t+1)=x(t)$ is a \emph{pause}. A \emph{horizon} $T\in\N$ fixes the times $[0,T)=\{0,\dots,T-1\}$. The objects
defined below depend only on the values of the walks at times below~$T$, so a walk of length~$T$ in the sense of the eighth
paper~\cite[Section~2.2]{EighthPaper2026} is a walk extended by $x(t)=x(T-1)$ for $t\ge T$. A walk is \emph{pause-free}
before~$T$ if $x(t+1)\ne x(t)$ whenever $t+1<T$. A tuple of $k$ walks $x_1,\dots,x_k$ is, in Lean, a map
$\Fin k\to(\N\to\Z)$; the walks may be coupled arbitrarily.

A \emph{long hop}, or simply a \emph{hop}, of a walk~$x$ before~$T$ is a pair $(p,u)$ of times with $p+2\le u<T$,
$x(u)=x(p)$ and $x(t)\ne x(p)$ for $p<t<u$ (\lean{BoundedNesting.IsHop}, \lean{BoundedNesting.hops}): two consecutive
visits to the \emph{cell} $z=x(p)$ that are not consecutive times. The hop \emph{spans} a time~$c$ if $p<c\le u$. Two hops
at one cell that span a common time coincide, and two hops with the same start coincide
(\lean{BoundedNesting.hop_unique}, \lean{BoundedNesting.hop_start_unique}); in particular a walk has at most $T$ hops
before~$T$. A hop is \emph{upper} if $x(t)>z$ for all $t\in(p,u)$ and \emph{lower} if $x(t)<z$ for all of them
(\lean{BoundedNesting.IsUpperHop}, \lean{BoundedNesting.IsLowerHop}). Its \emph{owner} is the first time $o\in(p,u)$ at which
the distance $\abs{x(o)-z}$ is largest (\lean{BoundedNesting.hopOwner}), and its \emph{extent} is this largest distance,
$\ext(p,u)=\max_{p<t<u}\abs{x(t)-z}$ (\lean{extent}).

\begin{lemma}[Sides and owners of hops]\label{lem:hops}
\formal{Walks.lean}{IsHop.upper_or_lower, isOwner_hopOwner_of_upper}
Let $x$ be a walk.
\begin{enumerate}[label=(\alph*)]
\item Every hop of $x$ is upper or lower.
\item The owner of an upper hop $(p,u)$ is the first time in $(p,u)$ at which $x$ is largest. (This part holds for every
$x\colon\N\to\Z$.)
\end{enumerate}
\end{lemma}

\begin{outline}
(a) The first step of a hop leaves its cell, so $x(p+1)>z$ or $x(p+1)<z$; a hop whose first step goes up lies strictly above
its cell, since the walk would otherwise return to~$z$ inside the hop (the eighth paper's \lean{BoundedNesting.above_of_step}),
and the case of a first step down is the same for $-x$. (b) Inside an upper hop $\abs{x(t)-z}=x(t)-z$, and the
characterization of \lean{BoundedNesting.hopOwner} by \lean{BoundedNesting.hopOwner_spec} gives the owner predicate
\lean{BoundedNesting.IsOwner}.
\end{outline}

The lower hops of~$x$ are the upper hops of~$-x$ (\lean{BoundedNesting.IsLowerHop.neg}), and the owner of a lower hop is the
first time inside it at which $x$ is smallest. The proof treats lower hops in this way, as upper hops of the reflected walk.
Upper hops do not cross: if $p<p'<u<u'$, then $x(p')>x(p)$ because $p'$ lies inside $(p,u)$, while $u$ lies inside $(p',u')$
and $x(u)=x(p)<x(p')$ (\lean{BoundedNesting.upper_no_cross}).

The \emph{walk DAG} of walks $x_1,\dots,x_k$ before~$T$ has the vertex set $[0,T)$, a \emph{state edge} $t\to t-1$ for
$1\le t<T$, and a \emph{hop edge} $u\to p$ for every pair $(p,u)$ that is a hop of at least one of the walks
(\lean{BoundedNesting.allHops}, \lean{BoundedNesting.hopPairs}); an edge is determined by its endpoints. A \emph{cut} is a set
of state edges, given by their upper ends $t\in[1,T)$, together with a set of hop edges. A path
$t_0\to t_1\to\dots\to t_n$ \emph{avoids} the cut if none of its edges belongs to it (\lean{BoundedNesting.AvoidPath}), and
the \emph{depth} of the cut is the largest number of hop edges on a path that avoids it (\lean{BoundedNesting.arcCount} counts
them). Paths run backwards in time.

\subsection{Cuts, labels, TCP and SR}\label{sec:prelim-tcp}

A \emph{label} with $q$ values is a map $\Psi\colon\N\to\{0,\dots,q-1\}$; only its values on $[0,T)$ matter. Its
\emph{descents} are the times $t\in[1,T)$ with $\Psi(t)<\Psi(t-1)$ (\lean{BoundedNesting.descentsOf}), and a hop edge
$u\to p$ is \emph{violated} if $\Psi(u)\le\Psi(p)$ (\lean{BoundedNesting.violatedOf}). The cut $F_\Psi$ consists of the state
edges at the descents and of the violated hop edges, and its size is the \emph{cost} of~$\Psi$ (\lean{labelCutSize}).

\begin{proposition}[Cuts are labels]\label{prop:labelcut}
\cited{\raggedright the eighth paper~\cite{EighthPaper2026}, Lemma~2.3; Lean-verified there
(\lean{BoundedNesting.labelCut_depth_walkDag}, from \lean{BoundedNesting.labelCut_depth} and
\lean{BoundedNesting.depth_le_of_label}).\par}
Let $\Psi$ be a label with $q\ge1$ values. For all walks and every horizon, the cut $F_\Psi$ of the walk DAG has depth at
most $q-1$.
\end{proposition}

Along a path that avoids $F_\Psi$, the label does not increase along a state edge and decreases strictly along a hop edge,
and it takes $q$ values. Every cut in this paper is of this form, so the task is to find labels with few values and small
cost.

\begin{definition}[TCP and SR]\label{def:tcp}
\cited{\raggedright the eighth paper~\cite{EighthPaper2026}, Definition~2.4 and Section~4.3, where these are Lean
definitions (\lean{BoundedNesting.TCPat}, \lean{BoundedNesting.TCP}, \lean{BoundedNesting.TCPnpAt},
\lean{BoundedNesting.TCPnp}, \lean{BoundedNesting.StackDag}, \lean{BoundedNesting.SRat}, \lean{BoundedNesting.SR}).\par}
Let $k,s\in\N$, $\eps\in\R$ and $D\in\N$.
\begin{enumerate}[label=(\alph*)]
\item $\TCP(k)$ holds at $(\eps,D)$ if for every horizon~$T$ and all walks $x_1,\dots,x_k$ the walk DAG before~$T$ has a cut
of at most $\eps T$ edges and depth at most~$D$. $\TCP_{\mathrm{np}}(k)$ holds at $(\eps,D)$ if the same is true for all
walks that are pause-free before~$T$.
\item $\TCP(k)$ is the statement that for every $\eps>0$ there is a~$D$ such that $\TCP(k)$ holds at $(\eps,D)$, and
likewise for $\TCP_{\mathrm{np}}(k)$.
\item An \emph{$s$-stack DAG} on $[0,N)$ has the \emph{spine edges} $t\to t-1$ and \emph{arcs} $u\to p$, given as pairs
$(p,u)$ with $p+2\le u<N$ and partitioned into $s$ families; within a family no two arcs cross, and every vertex is the left
endpoint of at most one arc of the family and the right endpoint of at most one. Cuts, paths and depth are as for walk DAGs,
with arcs in place of hop edges.
\item $\SR(s)$ is the statement that for every $\eps>0$ there is a~$D$ such that every $s$-stack DAG on $[0,N)$, for every
$N$, has a cut of at most $\eps N$ edges and depth at most~$D$.
\end{enumerate}
\end{definition}

\begin{proposition}[TCP is depth reduction for stack DAGs]\label{prop:stacksred}
\cited{\raggedright the eighth paper~\cite{EighthPaper2026}, Proposition~4.5; Lean-verified there
(\lean{BoundedNesting.walkDag_isStack}, \lean{BoundedNesting.srAt_of_tcpnpAt}, \lean{BoundedNesting.tcpAt_of_srAt},
\lean{BoundedNesting.sr_of_tcpnp}, \lean{BoundedNesting.tcp_of_sr}, \lean{BoundedNesting.tcp_of_tcpnp_two_mul}).\par}
\begin{enumerate}[label=(\alph*)]
\item The walk DAG of $k$ walks, pauses allowed, is a $2k$-stack DAG; its families are indexed by a walk and a side.
\item Every $s$-stack DAG on $[0,N)$ is realized by $s$ walks of length $2N$ without pauses, the arc $(p,u)$ of family~$i$
being the upper hop $(2p,2u)$ of walk~$i$, and a cut of their walk DAG pulls back to a cut of the stack DAG of no larger
size and depth.
\end{enumerate}
Consequently, if $\TCP_{\mathrm{np}}(s)$ holds at $(\eps/2,D)$, then every $s$-stack DAG on $[0,N)$ has a cut of at most
$\eps N$ edges and depth at most~$D$; if every $2k$-stack DAG on $[0,T)$ has a cut of at most $\eps T$ edges and depth at
most~$D$, then $\TCP(k)$ holds at $(\eps,D)$; and $\TCP_{\mathrm{np}}(2k)$ implies $\TCP(k)$.
\end{proposition}

We prove TCP for walks with pauses directly; the construction never uses pause-freeness. By \cref{prop:stacksred}, $\SR(s)$
then follows from the pause-free form for $s$ walks, a special case.

\begin{definition}[The main statement]\label{def:main}
\formal{Statement.lean}{MainStatement}
The \emph{main statement} is $\TCP(k)$ for every $k\in\N$.
\end{definition}

\begin{corollary}[Stack DAGs]\label{cor:main}
\formal{Statement.lean}{sr_of_main}
If the main statement holds, then $\SR(s)$ holds for every~$s$.
\end{corollary}

\begin{outline}
Pause-free walks are walks, so $\TCP(s)$ gives $\TCP_{\mathrm{np}}(s)$, and \cref{prop:stacksred} turns it into $\SR(s)$
(\lean{BoundedNesting.tcpnp_of_tcp}, \lean{BoundedNesting.sr_of_tcpnp}).
\end{outline}

The main statement is proved by the probabilistic method. For $k$ walks we build a finite family of labels, indexed by the
random choices, whose values lie below a bound $q$ that depends only on $k$ and~$\eps$ and whose average cost is at most
$\eps T$. Some member of the family is then a cut of the required kind. The following proposition is this argument; the rest
of the paper constructs the families and bounds their average cost.

\begin{proposition}[The probabilistic method for label cuts]\label{prop:avg}
\formal{Assembly.lean}{labelCutSize, exists_cut_of_avg}
Let $x_1,\dots,x_k$ be walks, $T$ a horizon, $\eps\in\R$, $\Omega$ a nonempty finite set, and $(\Psi_\omega)_{\omega\in\Omega}$
labels with values below~$q$. If $\sum_{\omega\in\Omega}\card{F_{\Psi_\omega}}\le\card\Omega\cdot\eps T$, then the walk DAG
before~$T$ has a cut of at most $\eps T$ edges and depth at most $q-1$.
\end{proposition}

\begin{outline}
Some $\omega\in\Omega$ has $\card{F_{\Psi_\omega}}\le\eps T$, since the sum is at most $\card\Omega$ times this bound
(Mathlib's \texttt{Finset.exists\_le\_of\_sum\_le}). The cut $F_{\Psi_\omega}$ consists of the descents and the violated hop
pairs of $\Psi_\omega$ (\lean{labelCutSize}), and \cref{prop:labelcut} bounds its depth by $q-1$ (\lean{exists_cut_of_avg}).
\end{outline}

\subsection{Owners and prominence}\label{sec:prelim-prom}

By \cref{lem:hops}(b), the owner of an upper hop is the first time inside it at which the walk is highest. We call the owners
of the upper hops of~$x$ its \emph{tops} (\lean{tops}), and the number of upper hops owned by a top~$\tau$ its \emph{height}
$H_\tau$ (\lean{height}); \cref{sec:ladders} is about the hops of one top. The proof measures the scale of a top by how far the
walk must descend from it before it reaches a higher point. Ties are broken in time: an earlier point of equal height counts
as higher, a later point only if it is strictly higher.

\begin{definition}[Prominence]\label{def:prom}
\formal{Prominence.lean}{colL, colR, peak, valley, prom, extent}
Let $x\colon\N\to\Z$, $T$ a horizon, $t<T$ and $v=x(t)$. Let $\min x$ denote the minimum of $x$ over $[0,T)$.
\begin{enumerate}[label=(\alph*)]
\item The \emph{left col} $L(t)$ is the minimum of $x$ over $[s_L,t]$, where $s_L$ is the last time $s<t$ with $x(s)\ge v$;
if there is no such time, $L(t)=\min x$.
\item The \emph{right col} $R(t)$ is the minimum of $x$ over $[t,s_R]$, where $s_R$ is the first time $s$ with $t<s<T$ and
$x(s)>v$; if there is no such time, $R(t)=\min x$.
\item The \emph{peak prominence} is $\max(0,\,v-\max(L(t),R(t)))$, the \emph{valley prominence} is the peak prominence of $-x$
at~$t$, and the \emph{prominence} $\prom(t)$ is the larger of the two.
\end{enumerate}
\end{definition}

A missing side contributes the global minimum, which is at most the other col, so it is ignored; if both sides are missing,
the peak prominence is $v-\min x$. For example, on the walk $0,1,2,1,2,3,2,1,0$ ($T=9$) the peak prominences are
$0,0,1,0,0,3,0,0,0$ and the valley prominences $3,0,0,1,0,0,0,0,3$; on $0,1,1,0$ the peak prominences are $0,1,0,0$, so of two
equal neighbouring maxima only the first is prominent. These values are checked by the kernel in \leanfile{Prominence.lean}.
By definition, the prominence does not change when $x$ is replaced by $-x$.

\begin{proposition}[Prominence]\label{prop:prom}
\formal{Prominence.lean}{sub_le_prom, sum_prom_le, sum_prom_le_two_mul, card_spanning_hops_lt}
\formal{TopProminence.lean}{height_le_prom, prom_le_of_left, prom_le_of_right}
Let $x\colon\N\to\Z$ and let $T$ be a horizon.
\begin{enumerate}[label=(P\arabic*)]
\item \emph{(Owners.)} If $(p,u)$ is an upper hop with $u<T$ and owner~$o$, then $\prom(o)\ge x(o)-x(p)$, the extent of the
hop.
\item \emph{(The prominence sum.)} If $x$ is a walk, then $\sum_{t<T}\prom(t)$ is at most the number of non-pause steps among
the first $T-1$ steps plus $\max x-\min x$ (over $[0,T)$), and hence at most $2(T-1)$.
\item \emph{(Spanning hops.)} For every time~$c$ and every $R\ge1$, fewer than $2R$ hops $(p,u)$ of~$x$ before~$T$ with
$p<c\le u$ have extent less than~$R$.
\item \emph{(Tops.)} If $x$ is a walk, every top $\tau$ has $\prom(\tau)\ge H_\tau$.
\item \emph{(Caps.)} Let $1\le t<T$ with $x(t-1)<x(t)$, and let $y\in\Z$. If some $s<t$ has $x(s)\ge x(t)$ and $x\ge y$
on $[s,t]$, or some $s$ with $t<s<T$ has $x(s)>x(t)$ and $x\ge y$ on $[t,s]$, then $\prom(t)\le x(t)-y$.
\end{enumerate}
\end{proposition}

\begin{outline}
(P1) The extent is attained at the owner (\lean{extent_le_owner}). The walk is below $x(o)$ on $[p,o)$, since $o$ is the
first maximum, and at most $x(o)$ on $(o,u]$. So the last earlier
time at least as high as~$o$ lies before~$p$, the first later time strictly higher lies after~$u$, and both cols are at most
$x(p)=x(u)$ (\lean{colL_le_of}, \lean{colR_le_of}, \lean{sub_le_peak}). (P2) A double count over levels (\lean{sum_peak_le}, \lean{sum_valley_le}). Fix a level~$y$ and call
a time~$\tau$ \emph{active} if $x(\tau)>y$ and both of its cols are at most~$y$, that is, if $y$ lies in
$[x(\tau)-\mathrm{peak}(\tau),x(\tau))$. Two active times never lie in one maximal run of times at which the walk is above~$y$:
of the two, the lower one, or the later one at equal height, has a point at least as high on the other one's side, which lifts
its col above~$y$ (\lean{active_unique}). Each such run starts at time~$0$ or at an up-step through~$y$
(\lean{card_active_le}), so summing over~$y$ bounds the peak prominences by the number of up-steps plus $x(0)-\min x$, and the
valley prominences by the number of down-steps plus $\max x-x(0)$ (\lean{sum_prom_le}). Both terms of the sum are at most
$T-1$ for a walk. (P3) The cell of such a hop is within $R-1$ of $x(c)$: either $c=u$ and $x(c)$ is the cell, or $c$ lies
inside the hop. These are $2R-1$ cells, and at each of them at most one hop spans~$c$
(\lean{BoundedNesting.hop_unique}). (P4) The hop of~$\tau$ at the lowest cell has extent $H_\tau$ (\cref{lem:ladder}), so
(P1) applies to it. (P5) Since the walk steps up into~$t$, the valley prominence of~$t$ vanishes
(\lean{valley_eq_zero_of_step_up}); the higher point~$s$, reached without going below~$y$, lifts the corresponding col above
$y-1$ (\lean{lt_colL}, \lean{lt_colR}).
\end{outline}

The owner bound (P1) makes every owner a turning point of positive prominence, which is how the label sees a hop
(\cref{lem:validity}); (P2) and (P3) bound the weight of turning points and the number of short hops at a time in the
accounting of \cref{sec:accounting}; (P4) and (P5) are the two estimates for tops that \cref{sec:ladders} needs.

\subsection{Bands and the random grid}\label{sec:prelim-grid}

Distances enter the label only through their binary orders of magnitude, grouped in a random grid of width~$h_0$.

\begin{definition}[Bands, groups and floors]\label{def:grid}
\formal{Basic.lean}{band, grp, flr}
The \emph{band} of $M\in\N$ is $\band(M)=\lfloor\log_2M\rfloor$ for $M\ge1$, and $\band(0)=0$. For a \emph{width} $h_0\ge1$
and an \emph{offset} $r\in\{0,\dots,h_0-1\}$, the \emph{group} of a band~$b$ is $G(b)=\lfloor(b+r)/h_0\rfloor$, and the
\emph{floor} of a group~$g$ is $\theta(g)=\max(h_0g-r,0)$. The \emph{grid floors} of the offset~$r$ are the floors $\theta(g)$
with $g\ge1$.
\end{definition}

The group map and the floors depend on $h_0$ and~$r$, which will be clear from the context. The groups are the blocks
$[\theta(g),\theta(g+1))$ of $h_0$ consecutive bands (the block of group~$0$ is shorter when $r>0$), and the offset shifts
the blocks. The offset will be uniform on $\{0,\dots,h_0-1\}$. Every probability bound that the paper needs about the grid
is then an upper bound on a number of offsets, and all of them come from one fact: an integer $\theta\ge1$ is a grid floor
for at most one offset. Indeed, $\theta=h_0g-r$ with $g\ge1$ means $h_0\mid\theta+r$, and two offsets below~$h_0$ that are
congruent modulo~$h_0$ are equal.

\begin{lemma}[The random grid]\label{lem:grid}
\formal{Grid.lean}{flr_grp_le, lt_flr_grp_succ, grp_mono, card_dvd_add_le_one, exists_floor_of_grp_lt, card_grp_ne_le, card_one_le_grp_le, card_flr_grp_succ_le, card_lt_flr_grp_le, sum_two_pow_flr_grp_le}
Let $h_0\ge1$, let the offset $r$ range over $\{0,\dots,h_0-1\}$, and let $a,b,c,s,\beta\in\N$. Subtraction of natural numbers
is truncated at~$0$.
\begin{enumerate}[label=(\alph*)]
\item $\theta(G(b))\le b<\theta(G(b)+1)$, and $G$ is nondecreasing.
\item For every $\theta\in\N$, at most one offset $r<h_0$ has $h_0\mid\theta+r$.
\item If $G(a)<G(b)$, then $h_0\mid\theta+r$ for some $\theta$ with $a<\theta\le b$.
\item If $a\le b$, then $G(a)\ne G(b)$ for at most $b-a$ offsets. Moreover $G(s)\ge1$ for at most $s$ offsets;
$\theta(G(b)+1)\le c$ for at most $c-b$ offsets; and $\theta(G(\beta))>s$ for at most $\beta-s$ offsets.
\item The sum of $2^{\theta(G(s))}$ over the offsets with $\theta(G(s))\ge1$ is at most $2^{s+1}$.
\end{enumerate}
\end{lemma}

\begin{outline}
(a) Division with remainder of $b+r$ by~$h_0$. (b) Two offsets $r,r'<h_0$ with $\theta+r\equiv\theta+r'\pmod{h_0}$ are equal
(\lean{card_dvd_add_le_one}). (c) The floor
$h_0G(b)-r$ lies in $(a,b]$. (d) Each count is bounded by the number of integers in an interval that are grid floors, each for
at most one offset by~(b). If $G(a)\ne G(b)$, then $G(a)<G(b)$ by~(a), and a grid floor lies in $(a,b]$ by~(c), an interval of
$b-a$ integers (\lean{card_grp_ne_le}). The other three reduce to this one: $G(s)\ge1$ means $G(0)\ne G(s)$, since $G(0)=0$;
$\theta(G(b)+1)\le c$ forces $G(b)<G(c)$; and $\theta(G(\beta))>s$ forces $G(s)<G(\beta)$. (e) Distinct offsets give distinct
positive floors $\theta(G(s))\le s$ (\lean{flr_grp_injOn}), so the sum is at most $\sum_{\theta\le s}2^\theta<2^{s+1}$.
\end{outline}

With $r$ uniform, part~(d) says that the groups of two bands $a\le b$ differ with probability at most $(b-a)/h_0$: this is how
band differences become costs in \cref{lem:cost}. Part~(e) prices the events of the accounting: for an event of scale~$s$, the
largest grid floor $\theta\le s$, weighted by $2^\theta$, averages at most $2^{s+1}/h_0$ over the offsets. The costs of one walk
are added up with the following exact sum, Lemma~E, which bounds the climb side of every ladder in \cref{sec:ladders}.

\begin{lemma}[Lemma E: band sums]\label{lem:bandsums}
\formal{BandSums.lean}{sum_band_sub_eq}
Let $N\in\N$ and $B=\band(N)$. Then $\sum_{x=1}^{N}(B-\band(x))=2^{B+1}-B-2$.
\end{lemma}

\begin{outline}
From $N$ to $N+1$ the band either stays or rises by one, and it rises only when $N+1$ is a power of two
(\lean{band_succ}); induction on~$N$ gives $\sum_{x=1}^{N}\band(x)+2^{B+1}=NB+B+2$ (\lean{sum_band_add_pow}), which is the
identity (\lean{sum_band_sub_eq}).
\end{outline}

Since $2^B\le N$ for $N\ge1$, the sum is at most $2N$: most integers up to~$N$ have a band close to $\band(N)$.
